\begin{mistake}{2026-04-19}{08}

\begin{problemcontent}
设 \(2\) 阶正交矩阵 \(A\) 的主对角线元素满足 \(a_{11} + 2 = a_{22}\)，则 \(A = \) \underline{\hspace{2cm}}。
\end{problemcontent}

\begin{solutioncontent}
设 \(2\) 阶正交矩阵 \(A = \begin{pmatrix} a_{11} & a_{12} \\ a_{21} & a_{22} \end{pmatrix}\)。

由正交矩阵的列向量为单位向量（模长为 \(1\)），可知矩阵中任意元素满足 \(|a_{ij}| \le 1\)，故有：
\[ -1 \le a_{11} \le 1, \quad -1 \le a_{22} \le 1 \]

将题设条件 \(a_{22} = a_{11} + 2\) 代入关于 \(a_{22}\) 的不等式：
\[ -1 \le a_{11} + 2 \le 1 \implies -3 \le a_{11} \le -1 \]

结合 \(-1 \le a_{11} \le 1\)，取交集可得唯一解：
\[ a_{11} = -1 \]
进而得到 \(a_{22} = a_{11} + 2 = 1\)。

再由列向量的单位性：
\[
\begin{aligned}
a_{11}^2 + a_{21}^2 &= 1 &\implies (-1)^2 + a_{21}^2 = 1 \implies a_{21} = 0 \\
a_{12}^2 + a_{22}^2 &= 1 &\implies a_{12}^2 + 1^2 = 1 \implies a_{12} = 0
\end{aligned}
\]

故该正交矩阵为：
\[ A = \begin{pmatrix} -1 & 0 \\ 0 & 1 \end{pmatrix} \]
\end{solutioncontent}

\mistakeknowledge{
\begin{itemize}
 \item \textbf{核心定理}：正交矩阵满足 \(A^TA=E\)，其列（行）向量均为标准正交基，任意元素必满足 \(a_{ij} \in [-1, 1]\)。
 \item \textbf{易错点}：设出四个未知数建立四次非线性方程组硬解，未第一时间利用元素的有界性去“夹逼”求解。
 \item \textbf{关联知识}：\(2\) 阶正交矩阵的标准参数形式仅有两种：旋转矩阵（主对角线元素相等）与反射矩阵（主对角线元素互为相反数）。
\end{itemize}}

\end{mistake}
